\chapter{Kết luận}
Luận văn này đã thực hiện khảo sát và đánh giá thực nghiệm nhiều thuật toán cắt tỉa mạng nơ-ron,
bao gồm GraSP, SynFlow, IMP, Early-Bird, Hybrid và Hybrid-Improve,
trên các cấu hình ResNet-20/CIFAR-10, ResNet-20/CIFAR-100 và ResNet-50/CIFAR-10.

Kết quả cho thấy các phương pháp hybrid đạt cân bằng tốt nhất giữa độ chính xác,
chi phí huấn luyện và tính ổn định.
Hybrid và Hybrid-Improve duy trì mức sụt giảm tương đối nhỏ ngay cả ở $90\%$ sparsity,
đồng thời ổn định hơn đáng kể so với các phương pháp khác.
Trong đó, Hybrid-Improve cho thấy độ ổn định tốt nhất giữa các lần chạy khác nhau.

GraSP nổi bật ở khả năng tạo phân phối độ thưa đồng đều và ổn định giữa các tầng,
trong khi SynFlow cho thấy độ nhạy cao với khởi tạo ngẫu nhiên.
IMP đạt hiệu suất tương đối tốt nhưng chi phí huấn luyện rất cao,
làm hạn chế tính thực tiễn ở quy mô lớn.
Early-Bird là phương pháp kém ổn định nhất tại mức độ thưa cực cao,
đặc biệt trên CIFAR-100.

Nhìn chung, kết quả thực nghiệm xác nhận rằng không tồn tại một thuật toán pruning tối ưu cho mọi trường hợp.
Hiệu quả của pruning phụ thuộc mạnh vào kiến trúc mạng, độ khó của bài toán,
mức độ thưa mục tiêu và năng lực tính toán của phần cứng.
Tuy nhiên, các chiến lược hybrid cho thấy tiềm năng lớn như một hướng tiếp cận cân bằng
giữa hiệu suất và chi phí, đặc biệt trong các bài toán yêu cầu mức nén cao nhưng vẫn cần duy trì độ chính xác.